\documentclass[12pt,a4paper]{article}
\usepackage[utf8]{inputenc}
\usepackage[T1]{fontenc}
\usepackage[brazil]{babel}
\usepackage{amsmath,amssymb,amsthm,mathtools}
\usepackage{graphicx}
\usepackage[margin=2.5cm,top=3cm,bottom=3cm]{geometry}
\usepackage{tikz}
\usetikzlibrary{arrows.meta,shapes.geometric,positioning,fit,backgrounds,calc,shadows.blur,decorations.pathmorphing,mindmap,trees,patterns,matrix}
\usepackage{pgfplots}
\pgfplotsset{compat=1.18}
\usepackage{fontawesome5}
\usepackage[ruled,vlined,linesnumbered]{algorithm2e}
\usepackage{listings}
\usepackage{xcolor}
\usepackage{booktabs}
\usepackage{multirow}
\usepackage{array}
\usepackage{tcolorbox}
\tcbuselibrary{breakable,skins,theorems}
\usepackage[hidelinks]{hyperref}
\usepackage{bookmark}
\setlength{\headheight}{14pt}
\usepackage{fancyhdr}
\usepackage{titlesec}
\usepackage{enumitem}
\usepackage{subcaption}
\usepackage{caption}
\definecolor{mineblue}{RGB}{0,90,180}
\definecolor{mineorange}{RGB}{230,115,0}
\definecolor{minegreen}{RGB}{34,139,34}
\definecolor{minered}{RGB}{180,20,40}
\definecolor{minegray}{RGB}{100,100,100}
\definecolor{codebg}{RGB}{250,250,250}
\newtcolorbox{definicaobox}[1]{title={\faIcon{book} #1},colback=blue!4!white,colframe=mineblue,fonttitle=\bfseries,breakable,enhanced}
\newtcolorbox{exemplobox}[1]{title={\faIcon{lightbulb} #1},colback=orange!4!white,colframe=mineorange,fonttitle=\bfseries,breakable,enhanced}
\newtcolorbox{dicabox}[1]{title={\faIcon{info-circle} #1},colback=green!4!white,colframe=minegreen,fonttitle=\bfseries,breakable,enhanced}
\newtcolorbox{atencaobox}[1]{title={\faIcon{exclamation-triangle} #1},colback=red!4!white,colframe=minered,fonttitle=\bfseries,breakable,enhanced}
\newtcolorbox{estadoartebox}[1]{title={\faIcon{rocket} #1},colback=purple!4!white,colframe=purple!60!black,fonttitle=\bfseries,breakable,enhanced}
\lstdefinestyle{mypython}{language=Python,backgroundcolor=\color{codebg},commentstyle=\color{minegray}\itshape,keywordstyle=\color{mineblue}\bfseries,stringstyle=\color{minegreen},basicstyle=\ttfamily\footnotesize,breaklines=true,frame=single,rulecolor=\color{minegray!40},numbers=left,numberstyle=\tiny\color{minegray},tabsize=4,showstringspaces=false}
\lstset{style=mypython}
\pagestyle{fancy}
\fancyhf{}
\fancyhead[L]{\small \textcolor{mineblue}{\textbf{Mineração de Dados}} | UEPA -- CCNT}
\fancyhead[R]{\small Módulo 8: Visualização de Dados}
\fancyfoot[C]{\small \thepage}
\renewcommand{\headrulewidth}{0.4pt}
\titleformat{\section}{\large\bfseries\color{mineblue}}{\thesection}{0.8em}{}[\titlerule]
\titleformat{\subsection}{\normalsize\bfseries\color{mineorange}}{\thesubsection}{0.8em}{}
\titleformat{\subsubsection}{\small\bfseries\color{minegray}}{\thesubsubsection}{0.8em}{}

\begin{document}

\begin{center}
  {\LARGE\bfseries\color{mineblue} Módulo 8}\\[0.4em]
  {\large\bfseries Técnicas de Visualização de Dados}\\[0.6em]
  {\small Mineração de Dados $\cdot$ 8.º Semestre $\cdot$ Engenharia de Software $\cdot$ UEPA/CCNT}\\[0.3em]
  {\small Carga horária do módulo: \textasciitilde{}10 horas}
\end{center}
\vspace{0.5em}
\hrule height 1.5pt
\vspace{1em}

% ============================================================
\section{Objetivos de Aprendizagem}
% ============================================================

Ao concluir este módulo, o estudante será capaz de:

\begin{enumerate}[label=\faIcon{check-circle}~\textbf{\arabic*.}, leftmargin=2cm, itemsep=0.3em]
  \item Compreender os princípios fundamentais de visualização de dados segundo Tufte e a Gramática dos Gráficos de Wilkinson.
  \item Selecionar o tipo de gráfico mais adequado para cada combinação de tipo de dado e objetivo comunicativo.
  \item Implementar gráficos estatísticos fundamentais (distribuição, relação, comparação, composição) com Matplotlib e Seaborn.
  \item Aplicar técnicas de redução de dimensionalidade (PCA, t-SNE, UMAP) para visualização de dados de alta dimensão.
  \item Criar visualizações interativas com Plotly e construir dashboards analíticos básicos com Dash e Streamlit.
  \item Visualizar resultados de modelos de mineração de dados: curvas ROC, matrizes de confusão, importância de atributos e gráficos SHAP.
  \item Identificar e evitar armadilhas comuns de visualização (eixos truncados, distorção 3D, excesso de informação).
  \item Aplicar boas práticas de acessibilidade em visualizações, incluindo paletas acessíveis a daltônicos e uso de texto alternativo.
\end{enumerate}

% ============================================================
\section{Princípios de Visualização de Dados}
% ============================================================

A visualização de dados é a representação gráfica de informação com o objetivo de comunicar padrões, tendências e anomalias de forma eficiente ao cérebro humano. Ela é simultaneamente uma ferramenta de \textbf{exploração} (EDA) e de \textbf{comunicação} (relatórios, dashboards).

\subsection{Por que Visualizar?}

\begin{exemplobox}{Quarteto de Anscombe (1973)}
Francis Anscombe demonstrou que quatro conjuntos de dados podem ter médias, variâncias, correlações e linhas de regressão idênticas, mas padrões completamente diferentes quando visualizados. Isso ilustra que estatísticas resumidas \textbf{nunca substituem} a visualização dos dados.

\begin{center}
\begin{tikzpicture}[scale=0.6]
  % Dataset I - linear
  \begin{scope}[xshift=0cm]
    \draw[->] (0,0)--(4,0); \draw[->] (0,0)--(0,4);
    \node[font=\tiny] at (2,-0.4) {Dataset I};
    \foreach \x/\y in {0.5/1.2,1.0/2.0,1.5/2.5,2.0/2.8,2.5/3.2,3.0/3.5,3.5/3.8}{
      \fill[mineblue] (\x,\y) circle(1.5pt);
    }
    \draw[mineorange,thick] (0,0.8)--(4,3.9);
  \end{scope}
  % Dataset II - curvilinear
  \begin{scope}[xshift=5cm]
    \draw[->] (0,0)--(4,0); \draw[->] (0,0)--(0,4);
    \node[font=\tiny] at (2,-0.4) {Dataset II};
    \foreach \x/\y in {0.3/0.6,0.8/1.8,1.3/2.7,1.8/3.2,2.3/3.4,2.8/3.2,3.3/2.8,3.8/1.8}{
      \fill[mineblue] (\x,\y) circle(1.5pt);
    }
    \draw[mineorange,thick] (0,0.9)--(4,3.8);
  \end{scope}
  % Dataset III - outlier
  \begin{scope}[xshift=10cm]
    \draw[->] (0,0)--(4,0); \draw[->] (0,0)--(0,4);
    \node[font=\tiny] at (2,-0.4) {Dataset III};
    \foreach \x/\y in {0.5/1.0,1.0/1.5,1.5/2.0,2.0/2.5,2.5/3.0,3.0/3.5,3.5/3.8}{
      \fill[mineblue] (\x,\y) circle(1.5pt);
    }
    \fill[minered] (1.8,3.9) circle(2pt);
    \draw[mineorange,thick] (0,0.8)--(4,3.7);
  \end{scope}
  % Dataset IV - vertical cluster
  \begin{scope}[xshift=15cm]
    \draw[->] (0,0)--(4,0); \draw[->] (0,0)--(0,4);
    \node[font=\tiny] at (2,-0.4) {Dataset IV};
    \foreach \y in {0.6,1.1,1.6,2.0,2.5,2.9,3.4}{
      \fill[mineblue] (2.0,\y) circle(1.5pt);
    }
    \fill[minered] (3.8,2.0) circle(2pt);
    \draw[mineorange,thick] (0,1.2)--(4,2.8);
  \end{scope}
\end{tikzpicture}
\end{center}
Todos têm: $\bar{x} \approx 9$, $\bar{y} \approx 7{,}5$, $r \approx 0{,}816$, mesma reta de regressão.
\end{exemplobox}

\textbf{Atributos pré-atentivos} são propriedades visuais detectadas pelo sistema nervoso humano antes do processamento consciente ($< 200$\,ms): cor, tamanho, orientação, posição, forma, movimento. Usá-los corretamente aumenta dramaticamente a eficiência da visualização.

\subsection{Princípios de Tufte (1983)}

Edward Tufte, no seminal \textit{The Visual Display of Quantitative Information}, estabelece princípios que guiam o design de visualizações eficientes:

\begin{itemize}
  \item \textbf{Data-ink ratio:} maximize a proporção de tinta usada para representar dados versus tinta decorativa. $\text{Data-ink ratio} = \frac{\text{tinta de dados}}{\text{tinta total}}$. Elimine grades desnecessárias, bordas, sombras.
  \item \textbf{Chartjunk:} evite elementos decorativos que não transmitem informação (padrões de preenchimento excessivos, imagens embutidas, perspectiva 3D desnecessária).
  \item \textbf{Lie Factor:} razão entre o efeito visual e o efeito real nos dados:
    \begin{equation}
      LF = \frac{\text{tamanho do efeito no gráfico}}{\text{tamanho do efeito nos dados}}
    \end{equation}
    $LF = 1$ é ideal; desvios sistemáticos configuram manipulação visual.
  \item \textbf{Small multiples:} conjuntos de gráficos com a mesma escala e estrutura para facilitar comparação temporal ou entre grupos.
  \item \textbf{Densidade de dados:} maximize a quantidade de informação por unidade de área.
\end{itemize}

\subsection{Grammar of Graphics (Wilkinson, 2005)}

Leland Wilkinson formalizou a \textbf{Gramática dos Gráficos} como um sistema composicional com sete componentes:

\begin{enumerate}
  \item \textbf{Data:} conjunto de dados subjacente.
  \item \textbf{Aesthetics (aes):} mapeamento de variáveis a propriedades visuais (x, y, cor, tamanho, forma, opacidade).
  \item \textbf{Geometry (geom):} tipo de elemento geométrico (ponto, linha, barra, polígono).
  \item \textbf{Facets:} subdivisão em painéis por variável categórica.
  \item \textbf{Statistics (stat):} transformações estatísticas (binning, suavização, contagem).
  \item \textbf{Coordinates (coord):} sistema de coordenadas (cartesiano, polar, mapa).
  \item \textbf{Themes:} aparência visual não-data (fonte, cor de fundo, grade).
\end{enumerate}

Implementações: \textbf{ggplot2} (R, Wickham 2010), \textbf{plotnine} (Python), \textbf{Plotly} (implementa parcialmente a gramática via \texttt{go.Figure}).

\subsection{Escolha do Tipo de Gráfico}

\begin{center}
\begin{tikzpicture}[
  question/.style={diamond,draw=mineblue,fill=mineblue!10,minimum width=2.8cm,minimum height=1cm,align=center,font=\tiny,aspect=2.5},
  answer/.style={rectangle,rounded corners,draw=minegreen,fill=minegreen!10,minimum width=2.5cm,minimum height=0.7cm,align=center,font=\tiny},
  arr/.style={->,thick,minegray}]

\node[question] (q1) at (0,0) {O que quer\\mostrar?};
\node[question] (q2) at (-4.5,-2.5) {Número de\\variáveis?};
\node[question] (q3) at (4.5,-2.5) {Tipo de\\comparação?};

\node[answer] (a1) at (-6.5,-5) {Histograma\\Boxplot / Violino};
\node[answer] (a2) at (-2.5,-5) {Scatter / Hexbin\\Heatmap de Corr.};
\node[answer] (a3) at (2.5,-5) {Barras / Linhas\\Radar / Strip};
\node[answer] (a4) at (6.5,-5) {Treemap\\Sunburst / Waffle};

\draw[arr] (q1)--(-3.5,0)--(-4.5,-1.8) node[pos=0.3,above,font=\tiny]{Distribuição};
\draw[arr] (q1)--(3.5,0)--(4.5,-1.8) node[pos=0.3,above,font=\tiny]{Comparação};
\draw[arr] (q2)--(-6.5,-4.3) node[midway,left,font=\tiny]{1 variável};
\draw[arr] (q2)--(-2.5,-4.3) node[midway,right,font=\tiny]{2+ variáveis};
\draw[arr] (q3)--(2.5,-4.3) node[midway,left,font=\tiny]{Temporal};
\draw[arr] (q3)--(6.5,-4.3) node[midway,right,font=\tiny]{Composição};
\end{tikzpicture}
\end{center}

\begin{table}[ht]
\centering
\small
\caption{Guia de seleção de tipos de gráfico}
\begin{tabular}{@{}p{2.8cm}p{2.2cm}p{3cm}p{2.5cm}p{2.5cm}@{}}
\toprule
\textbf{Tipo de Gráfico} & \textbf{Tipo de Dado} & \textbf{Caso de Uso} & \textbf{Exemplo} & \textbf{Ferramentas} \\
\midrule
Histograma & Numérico & Distribuição univariada & Distribuição de idades & matplotlib, seaborn \\
Boxplot & Num.\ por categ. & Comparar distribuições & Notas por turma & seaborn, plotly \\
Scatter Plot & 2 Numéricos & Relação/correlação & Renda vs consumo & matplotlib, plotly \\
Heatmap & Matriz & Correlações, matrizes & Corr.\ de features & seaborn, plotly \\
Barras & Categ.\ + Num. & Comparação de grupos & Vendas por região & matplotlib, plotly \\
Linha & Temporal & Séries temporais & Acurácia por época & matplotlib, plotly \\
Treemap & Hierárquico & Composição aninhada & Receita por produto & plotly, squarify \\
Violino & Num.\ por categ. & Distribuição + densidade & Salário por cargo & seaborn \\
Choropleth & Geoespacial & Distribuição geográfica & IDH por estado & plotly, folium \\
Parallel Coord. & Multivariado & Alta dimensão & Comparar clusters & plotly \\
\bottomrule
\end{tabular}
\end{table}

% ============================================================
\section{Gráficos Estatísticos Fundamentais}
% ============================================================

\subsection{Distribuição}

\begin{itemize}
  \item \textbf{Histograma:} divide dados em bins e conta frequência. Use para exploração inicial e detecção de bimodalidade. Cuidado com número de bins.
  \item \textbf{KDE (Kernel Density Estimate):} versão suavizada do histograma, sem a arbitrariedade do binning.
  \item \textbf{Boxplot:} mediana, IQR, whiskers ($\pm 1{,}5 \times IQR$) e outliers. Excelente para comparação de grupos.
  \item \textbf{Violin Plot:} combina boxplot com KDE; mostra a distribuição completa, especialmente útil quando bimodal.
  \item \textbf{Q-Q Plot:} compara a distribuição observada com uma teórica (e.g., normal); pontos na diagonal indicam ajuste.
\end{itemize}

\begin{lstlisting}[caption={Gráficos de distribuição com seaborn}]
import seaborn as sns, matplotlib.pyplot as plt

fig, axes = plt.subplots(1, 3, figsize=(14, 4))
sns.histplot(data=df, x='feature', kde=True, ax=axes[0])
axes[0].set_title('Histograma + KDE')

sns.boxplot(data=df, x='category', y='feature', ax=axes[1])
axes[1].set_title('Boxplot por Grupo')

sns.violinplot(data=df, x='category', y='feature',
               inner='quartile', ax=axes[2])
axes[2].set_title('Violin Plot')
plt.tight_layout(); plt.show()
\end{lstlisting}

\subsection{Relação entre Variáveis}

\begin{itemize}
  \item \textbf{Scatter Plot:} relação entre duas variáveis contínuas. Adicione cor/tamanho para terceira/quarta dimensão.
  \item \textbf{Bubble Chart:} scatter com tamanho proporcional a uma terceira variável quantitativa.
  \item \textbf{Hexbin:} versão do scatter para grandes datasets (evita overplotting com contagem por célula hexagonal).
  \item \textbf{Joint Plot:} scatter central + distribuições marginais. Seaborn \texttt{JointGrid}.
  \item \textbf{Correlation Heatmap:} matriz de correlações com codificação de cor; use máscara para triângulo superior.
  \item \textbf{Pair Plot (SPLOM):} todos os scatter plots par a par em uma grade. Ideal para EDA inicial.
\end{itemize}

\subsection{Comparação}

\begin{itemize}
  \item \textbf{Gráfico de barras agrupado:} compara múltiplos grupos por categoria. Evite mais de 5 grupos.
  \item \textbf{Gráfico de barras empilhado:} mostra contribuição de subgrupos ao total.
  \item \textbf{Lollipop chart:} alternativa mais limpa ao gráfico de barras, alta data-ink ratio.
  \item \textbf{Radar / Spider Chart:} compara perfis multivariados em múltiplos eixos radiais; bom para comparação de clusters.
  \item \textbf{Time Series Line Plot:} tendências temporais com múltiplas séries. Adicione bandas de confiança.
  \item \textbf{Strip / Swarm Plot:} mostra todos os pontos individuais, complementa o boxplot.
\end{itemize}

\subsection{Composição}

\begin{itemize}
  \item \textbf{Gráfico de pizza:} use com moderação -- apenas com 2--4 categorias, quando a proporção relativa ao todo é o foco. \textbf{Nunca use 3D.}
  \item \textbf{Treemap:} área proporcional ao valor; excelente para dados hierárquicos.
  \item \textbf{Sunburst:} versão radial do treemap para hierarquias multinível.
  \item \textbf{Waffle Chart:} grade de quadrados onde cada quadrado representa percentual; mais preciso visualmente que pizza.
\end{itemize}

\subsection{Geoespacial}

\begin{itemize}
  \item \textbf{Choropleth Maps:} polígonos geográficos coloridos por valor (e.g., IDH por estado). Implementação com Plotly ou Folium.
  \item \textbf{Scatter Geo:} pontos em mapa para dados com coordenadas (lat/lon).
  \item \textbf{Folium/Leaflet:} mapas interativos no browser com layers, clusters e tooltips.
  \item \textbf{Cartopy:} projeções cartográficas científicas integradas ao Matplotlib.
\end{itemize}

% ============================================================
\section{Visualização de Alta Dimensão}
% ============================================================

Datasets de mineração de dados frequentemente possuem dezenas ou centenas de features. A visualização direta é impossível em $d > 3$, tornando necessárias técnicas de redução de dimensionalidade para projeção em 2D ou 3D.

\subsection{Redução para 2D/3D}

\subsubsection{PCA (Principal Component Analysis)}

A PCA projeta os dados nas direções de máxima variância:
\begin{equation}
  \mathbf{Z} = \mathbf{X} \mathbf{W}, \quad \text{onde } \mathbf{W} = \text{eigenvectors de } \mathbf{X}^T\mathbf{X}
\end{equation}

O \textbf{PCA biplot} mostra simultaneamente as observações (pontos) e os loadings das variáveis originais (vetores), revelando quais variáveis mais contribuem para cada componente principal.

\subsubsection{t-SNE (van der Maaten \& Hinton, 2008)}

t-SNE minimiza a divergência KL entre a distribuição de pares no espaço original (Gaussiana) e no espaço reduzido (t-Student):
\begin{equation}
  KL(P \| Q) = \sum_i \sum_j p_{ij} \log \frac{p_{ij}}{q_{ij}}
\end{equation}

\textbf{Características:} não-linear, estocástico, excelente para visualizar clusters. \textbf{Parâmetro crítico:} perplexidade (5--50). \textbf{Limitação:} lento para $n > 10.000$, não preserva estrutura global.

\subsubsection{UMAP (McInnes et al., 2018)}

UMAP baseia-se em geometria Riemanniana e teoria de topologia fuzzy. Preserva tanto estrutura local quanto global, é mais rápido que t-SNE e permite transformar novos pontos (fit/transform).

\begin{table}[ht]
\centering
\caption{Comparação: PCA vs t-SNE vs UMAP}
\begin{tabular}{@{}lllll@{}}
\toprule
\textbf{Critério} & \textbf{PCA} & \textbf{t-SNE} & \textbf{UMAP} \\
\midrule
Tipo & Linear & Não-linear & Não-linear \\
Velocidade & Muito rápido & Lento ($O(n^2)$) & Rápido ($O(n \log n)$) \\
Estrutura global & Boa & Ruim & Boa \\
Estrutura local & Moderada & Excelente & Excelente \\
Determinismo & Sim & Não & Não (seed fix) \\
Parâmetros & Nenhum & Perplexidade & n\_neighbors, min\_dist \\
Transform novos & Sim & Não & Sim \\
\bottomrule
\end{tabular}
\end{table}

\begin{lstlisting}[caption={Comparação PCA, t-SNE, UMAP para visualização}]
from sklearn.manifold import TSNE
from sklearn.decomposition import PCA
import umap

# PCA
pca = PCA(n_components=2)
X_pca = pca.fit_transform(X_scaled)

# t-SNE
tsne = TSNE(n_components=2, perplexity=30, random_state=42)
X_tsne = tsne.fit_transform(X_scaled)

# UMAP
reducer = umap.UMAP(n_neighbors=15, min_dist=0.1, random_state=42)
X_umap = reducer.fit_transform(X_scaled)
\end{lstlisting}

\subsection{Parallel Coordinates}

Cada variável é representada como um eixo vertical paralelo, e cada observação é uma linha poligonal que cruza todos os eixos. Permite identificar clusters, outliers e correlações em dados de alta dimensão. Plotly oferece \texttt{go.Parcoords} com interatividade (brushing por eixo).

\subsection{Matriz de Dispersão (SPLOM)}

Todos os scatter plots pairwise em uma grade $d \times d$. Diagonal pode mostrar histogramas ou KDE. Codificação por cor revela separação de classes. \texttt{plotly.express.scatter\_matrix()} gera SPLOM interativo com hover e seleção.

% ============================================================
\section{Visualização Interativa}
% ============================================================

\subsection{Plotly}

Plotly é a biblioteca de visualização interativa mais popular para Python, gerando gráficos baseados em D3.js:

\begin{itemize}
  \item \textbf{Plotly Express:} API de alto nível para gráficos com uma linha de código. \texttt{px.scatter(), px.bar(), px.choropleth()}.
  \item \textbf{go.Figure:} API de baixo nível para layouts customizados, múltiplos eixos, anotações.
  \item \textbf{Animações:} parâmetro \texttt{animation\_frame} em Express ou \texttt{frames} em Figure.
  \item \textbf{Subplots:} \texttt{make\_subplots(rows, cols)} com compartilhamento de eixos.
\end{itemize}

\begin{lstlisting}[caption={Dashboard interativo com Plotly subplots}]
import plotly.graph_objects as go
from plotly.subplots import make_subplots

fig = make_subplots(rows=2, cols=2,
    subplot_titles=["ROC Curve", "Confusion Matrix",
                    "Feature Importance", "Learning Curve"])

fig.add_trace(go.Scatter(x=fpr, y=tpr, name=f"AUC={auc:.2f}"),
              row=1, col=1)
fig.add_trace(go.Heatmap(z=cm, colorscale="Blues"), row=1, col=2)
fig.add_trace(go.Bar(x=importances, y=features, orientation='h'),
              row=2, col=1)
fig.update_layout(title="Dashboard de Avaliacao de Modelo",
                  height=700)
fig.show()
\end{lstlisting}

\subsection{Bokeh}

Bokeh oferece interatividade server-side avançada:
\begin{itemize}
  \item \textbf{Linked plots (brushing):} seleção em um gráfico reflete em todos os outros conectados.
  \item \textbf{Bokeh Server:} callbacks Python ativados por interações do usuário.
  \item \textbf{Hover tools:} tooltips customizados com HTML.
\end{itemize}

\subsection{Dash}

Dash (Plotly) é um framework para construção de aplicações analíticas web em Python puro:

\begin{lstlisting}[caption={Aplicação Dash básica com callback}]
import dash
from dash import dcc, html, Input, Output
import plotly.express as px

app = dash.Dash(__name__)
app.layout = html.Div([
    dcc.Dropdown(id='x-axis', options=['feature1','feature2'], value='feature1'),
    dcc.Graph(id='scatter-plot')
])

@app.callback(Output('scatter-plot', 'figure'),
              Input('x-axis', 'value'))
def update_figure(x_col):
    return px.scatter(df, x=x_col, y='target', color='cluster')

if __name__ == '__main__':
    app.run(debug=True)
\end{lstlisting}

\subsection{Streamlit}

Streamlit permite criar dashboards analíticos em minutos:
\begin{lstlisting}[caption={Dashboard Streamlit para resultados de DM}]
import streamlit as st
import plotly.express as px

st.title("Dashboard: Mineracao de Dados - UEPA")
with st.sidebar:
    k = st.slider("Numero de clusters", 2, 10, 4)

tab1, tab2 = st.tabs(["Clustering", "Classificacao"])
with tab1:
    st.plotly_chart(px.scatter(df_tsne, x='x', y='y', color=f'cluster_{k}'))
with tab2:
    st.dataframe(df_metrics, use_container_width=True)
\end{lstlisting}

% ============================================================
\section{Visualização de Resultados de Mineração de Dados}
% ============================================================

\subsection{Avaliação de Modelos}

\begin{itemize}
  \item \textbf{Curva ROC:} True Positive Rate vs False Positive Rate para diferentes thresholds. AUC-ROC mede a área sob a curva; AUC $= 1$ é perfeito, $0{,}5$ é aleatório.
  \item \textbf{Curva Precision-Recall:} preferível à ROC em datasets muito desbalanceados. AUC-PR é mais informativa quando a classe positiva é rara.
  \item \textbf{Matriz de Confusão (Heatmap):} visualização de TP, FP, FN, TN. Normalize por linha para mostrar taxas de acerto por classe.
  \item \textbf{Curvas de Aprendizagem:} desempenho no treino e validação em função do tamanho do dataset. Diagnostica underfitting (ambas baixas) e overfitting (treino alta, validação baixa).
\end{itemize}

\subsection{Interpretabilidade Visual}

\subsubsection{SHAP (SHapley Additive exPlanations)}

SHAP (Lundberg \& Lee, 2017) atribui a cada feature a contribuição marginal para a predição individual:

\begin{itemize}
  \item \textbf{Summary Plot (Beeswarm):} todos os pontos de todos os exemplos; cor = valor da feature, eixo x = impacto SHAP. Revela distribuição de importância.
  \item \textbf{Bar Summary:} média do valor absoluto SHAP por feature; ranking global.
  \item \textbf{Waterfall Plot:} para uma predição individual; detalha contribuição de cada feature a partir do valor base.
  \item \textbf{Dependence Plot:} SHAP value vs valor da feature, colorido por feature de interação.
\end{itemize}

\begin{lstlisting}[caption={Visualizações SHAP para interpretabilidade}]
import shap

explainer = shap.TreeExplainer(model)
shap_values = explainer.shap_values(X_test)

# Summary beeswarm
shap.summary_plot(shap_values, X_test, feature_names=feature_names)

# Waterfall para observacao 0
shap.waterfall_plot(shap.Explanation(
    values=shap_values[0],
    base_values=explainer.expected_value,
    data=X_test.iloc[0],
    feature_names=feature_names))
\end{lstlisting}

\subsubsection{LIME}

LIME (Ribeiro et al., 2016) explica predições individuais aproximando o modelo com um modelo linear local. A visualização mostra features com contribuição positiva (verde) e negativa (vermelho) para a predição.

\subsection{Clustering Visualization}

\begin{itemize}
  \item \textbf{t-SNE / UMAP colorido por cluster:} visualização 2D dos clusters descobertos.
  \item \textbf{Dendrograma:} hierarquia de fusões no clustering aglomerativo.
  \item \textbf{Silhouette Plot:} gráfico de coeficientes de silhueta por amostra e cluster. Coeficiente $s_i \in [-1, 1]$; $s_i > 0$ bem clustered, $s_i < 0$ mal clustered.
  \item \textbf{Elbow Curve:} inércia vs $k$; o ``cotovelo'' sugere o $k$ ótimo.
  \item \textbf{Radar Chart dos Centroides:} perfil multivariado de cada cluster em gráfico radial.
\end{itemize}

\subsection{Associação}

\begin{itemize}
  \item \textbf{Grafo de regras:} nós são itemsets, arestas são regras; espessura = confiança, cor = lift.
  \item \textbf{Heatmap de co-ocorrência:} frequência de pares de itens.
  \item \textbf{Bubble Chart:} suporte (eixo x) vs confiança (eixo y), tamanho da bolha = lift. Permite identificar regras interessantes visualmente.
\end{itemize}

% ============================================================
\section{Dashboards Analíticos}
% ============================================================

\begin{definicaobox}{Dashboard Analítico}
Um dashboard é uma interface visual que consolida múltiplos indicadores e gráficos em uma única tela, permitindo monitoramento de KPIs e análise exploratória em tempo real.
\end{definicaobox}

\textbf{Princípios de design de dashboards:}
\begin{enumerate}
  \item \textbf{Hierarquia visual:} os itens mais importantes devem ser maiores e mais proeminentes (KPI cards no topo).
  \item \textbf{Alinhamento:} alinhe elementos em uma grade para criar ordem visual.
  \item \textbf{Proximidade:} agrupe elementos relacionados visualmente.
  \item \textbf{Contraste:} diferencie seções com cor ou espaço.
  \item \textbf{Drill-down:} permita ao usuário navegar do macro para o micro.
  \item \textbf{Mobile responsiveness:} layouts que se adaptam a telas menores.
\end{enumerate}

\paragraph{Mockup de Dashboard Analítico para DM}

\begin{center}
\begin{tikzpicture}[scale=0.7]
% Header
\fill[mineblue] (0,9) rectangle (16,10);
\node[text=white,font=\bfseries] at (8,9.5) {Dashboard: Mineração de Dados -- UEPA};

% KPI cards (top row)
\foreach \x/\label/\val/\col in {0.3/Acurácia/92\%/minegreen, 4.3/Precisão/89\%/mineblue, 8.3/Recall/91\%/mineorange, 12.3/F1-Score/90\%/purple}{
  \fill[\col!20,draw=\col] (\x,7.2) rectangle (\x+3.4,8.8);
  \node[font=\small\bfseries] at (\x+1.7,8.3) {\val};
  \node[font=\tiny] at (\x+1.7,7.6) {\label};
}

% Left chart
\fill[minegray!10,draw=minegray!40] (0.3,3.5) rectangle (7.8,7.0);
\node[font=\small,minegray] at (4,5.2) {Importância de Atributos};
% Bar chart sketch
\foreach \bx/\bh in {1.0/2.0, 1.8/1.4, 2.6/1.8, 3.4/1.1, 4.2/0.8}{
  \fill[mineblue!50] (\bx,3.8) rectangle (\bx+0.6,3.8+\bh*0.9);
}

% Right chart
\fill[minegray!10,draw=minegray!40] (8.2,3.5) rectangle (15.7,7.0);
\node[font=\small,minegray] at (12,5.2) {Curva ROC};
\draw[mineblue,thick] (9.0,3.8) .. controls (10,5.5) and (11,6.2) .. (15.0,6.7);
\draw[minegray,dashed] (9.0,3.8)--(15.0,6.7);

% Bottom charts
\fill[minegray!10,draw=minegray!40] (0.3,0.3) rectangle (5.0,3.2);
\node[font=\small,minegray] at (2.65,1.7) {Distribuição};

\fill[minegray!10,draw=minegray!40] (5.4,0.3) rectangle (10.1,3.2);
\node[font=\small,minegray] at (7.75,1.7) {Matriz de Confusão};

\fill[minegray!10,draw=minegray!40] (10.5,0.3) rectangle (15.7,3.2);
\node[font=\small,minegray] at (13.1,1.7) {Clustering};
\end{tikzpicture}
\end{center}

% ============================================================
\section{Boas Práticas e Armadilhas Comuns}
% ============================================================

\subsection{Erros Comuns}

\begin{atencaobox}{Principais Armadilhas na Visualização}
\begin{itemize}
  \item \textbf{Eixo y truncado:} cortar o eixo y para exagerar diferenças pequenas. Lie Factor elevado. Sempre comece o eixo y em zero para gráficos de barra.
  \item \textbf{Gráfico de pizza 3D:} perspectiva deforma as áreas, impossibilitando comparação precisa.
  \item \textbf{Excesso de cores:} mais de 7 categorias em uma legenda de cor torna o gráfico ilegível. Agrupe categorias menores em ``Outros''.
  \item \textbf{Ausência de unidades e rótulos:} eixos sem título e unidade tornam o gráfico ininterpretável.
  \item \textbf{Gráficos de duplo eixo:} dois eixos y diferentes permitem manipulação visual da correlação aparente entre séries.
  \item \textbf{Overplotting:} em scatter com muitos pontos, use transparência (\texttt{alpha}) ou hexbin.
  \item \textbf{Rainbow colormap:} \textit{jet}/\textit{rainbow} não são perceptualmente uniformes; use \textit{viridis} ou \textit{plasma}.
\end{itemize}
\end{atencaobox}

\subsection{Acessibilidade}

\begin{itemize}
  \item \textbf{Daltonismo:} cerca de 8\% dos homens têm dificuldade de distinguir vermelho e verde. Use paletas acessíveis: \textbf{ColorBrewer}, \textbf{viridis}, \textbf{cividis}, \textbf{Wong} (8 cores acessíveis).
  \item \textbf{Alto contraste:} razão mínima de 4,5:1 entre texto e fundo (WCAG 2.1 nível AA).
  \item \textbf{Alt text:} para imagens em relatórios web/PDF, descreva o padrão principal do gráfico.
  \item \textbf{Não dependência de cor exclusiva:} use combinação de cor + forma + padrão para codificar categorias.
\end{itemize}

% ============================================================
\section{Estado da Arte (2020--2024)}
% ============================================================

\begin{estadoartebox}{Visualização Exploratória Automática (Auto-EDA)}
Ferramentas de EDA automatizada eliminam a necessidade de escrever código para a exploração inicial de datasets. \textbf{ydata-profiling} (ex-Pandas Profiling) gera relatórios HTML completos com estatísticas, correlações, valores ausentes e alertas de qualidade. \textbf{Sweetviz} compara dois datasets (treino vs teste) destacando diferenças de distribuição. \textbf{AutoViz} gera automaticamente os gráficos mais relevantes. \textbf{lux} recomenda visualizações contextuais no Jupyter Notebook. \textbf{D-Tale} oferece exploração interativa com manipulação point-and-click.
\end{estadoartebox}

\begin{estadoartebox}{Visualização Explicável (XAI Viz)}
SHAP tornou-se o padrão da indústria para explicabilidade de modelos de ML, com plots integrados ao scikit-learn e XGBoost. \textbf{DiCE} (Diverse Counterfactual Explanations) gera exemplos contrafactuais visuais: ``o que precisaria mudar para inverter a predição?''. O \textbf{What-if Tool} do Google permite exploração visual interativa de modelos no TensorBoard. \textbf{Alibi} oferece explicações anchor e contrafactuais. A tendência é integrar XAI diretamente nos dashboards de produção para auditoria contínua.
\end{estadoartebox}

\begin{estadoartebox}{BI Moderno e DataViz}
\textbf{Apache Superset} (open-source) compete com Tableau e Power BI para dashboards corporativos. \textbf{Metabase} simplifica o acesso de usuários não técnicos a dados. \textbf{Grafana} tornou-se padrão para dados de séries temporais e monitoramento de sistemas. \textbf{Observable Plot} (JavaScript, 2022) é a evolução do D3.js, mais acessível. \textbf{Evidence} e \textbf{Rill} emergem como ferramentas de BI baseadas em SQL com output estático. Vibe-coded dashboards com LLMs (Text-to-Viz) começam a democratizar a criação de visualizações para usuários não programadores.
\end{estadoartebox}

% ============================================================
\section{Exercícios}
% ============================================================

\begin{enumerate}[label=\textbf{\arabic*.}, itemsep=0.5em]
  \item \textbf{[Quarteto de Anscombe]} Crie manualmente os quatro datasets do Quarteto de Anscombe usando NumPy. Calcule e imprima as estatísticas de cada um (média, variância, correlação, coeficientes de regressão linear). Em seguida, plote todos em um grid $2 \times 2$ com matplotlib. O que esse exercício ensina sobre a necessidade de visualização?

  \item \textbf{[Design de Dashboard]} Para o projeto de DM do módulo anterior (classificação ou clustering), projete um dashboard analítico no papel (wireframe) com: (a) 4 KPI cards no topo, (b) gráfico de avaliação do modelo, (c) visualização dos dados com PCA/t-SNE, (d) importância de atributos. Implemente o wireframe como figura Plotly \texttt{make\_subplots}.

  \item \textbf{[Plotly Interativo]} Crie um scatter plot interativo com Plotly Express usando o dataset Iris: (a) cores por espécie, (b) tamanho dos pontos proporcional ao comprimento da pétala, (c) hover customizado mostrando todas as 4 features, (d) botões de animação para alternar entre pares de features.

  \item \textbf{[Visualização SHAP]} Treine um XGBoost no dataset de câncer de mama (sklearn). Gere: (a) beeswarm summary plot, (b) bar summary plot, (c) waterfall plot para 3 exemplos (um verdadeiro positivo, um falso positivo, um verdadeiro negativo). Interprete cada plot em 3 frases.

  \item \textbf{[t-SNE vs UMAP]} Carregue o dataset MNIST (use \texttt{fetch\_openml('mnist\_784')} ou uma amostra de 5000 pontos). Aplique PCA (primeiro reduzir a 50D), depois t-SNE e UMAP. Compare visualmente a separação das 10 classes. Meça o tempo de execução de cada método. Qual preserva melhor a estrutura global?

  \item \textbf{[Auto-EDA Report]} Carregue qualquer dataset com pelo menos 10 colunas e 500 linhas. Execute \texttt{ydata\_profiling.ProfileReport} e salve o relatório HTML. Identifique: (a) 3 problemas de qualidade de dados detectados automaticamente, (b) 2 correlações interessantes entre variáveis, (c) distribu\-ições problemáticas (multimodal, skewed, outliers).

  \item \textbf{[Acessibilidade]} Crie o mesmo gráfico de barras em três versões: (a) versão padrão com cores arbitrárias, (b) versão com paleta ColorBrewer acessível a daltônicos, (c) versão com hachuras (patterns) além de cores. Use o simulador de daltonismo Coblis ou matplotlib.colors para verificar a acessibilidade. Documente as diferenças.
\end{enumerate}

% ============================================================
\section{Referências}
% ============================================================

\begin{thebibliography}{99}

\bibitem{tufte1983} Tufte, E. R. (1983). \textit{The Visual Display of Quantitative Information}. Graphics Press.

\bibitem{wilkinson2005} Wilkinson, L. (2005). \textit{The Grammar of Graphics} (2nd ed.). Springer.

\bibitem{vandermaaten2008} van der Maaten, L., \& Hinton, G. (2008). Visualizing data using t-SNE. \textit{Journal of Machine Learning Research}, 9, 2579--2605.

\bibitem{mcinnes2018} McInnes, L., Healy, J., \& Melville, J. (2018). UMAP: Uniform Manifold Approximation and Projection for Dimension Reduction. \textit{arXiv:1802.03426}.

\bibitem{wickham2010} Wickham, H. (2010). A layered grammar of graphics. \textit{Journal of Computational and Graphical Statistics}, 19(1), 3--28.

\bibitem{cairo2016} Cairo, A. (2016). \textit{The Truthful Art: Data, Charts, and Maps for Communication}. New Riders.

\bibitem{munzner2014} Munzner, T. (2014). \textit{Visualization Analysis and Design}. CRC Press.

\bibitem{few2012} Few, S. (2012). \textit{Show Me the Numbers: Designing Tables and Graphs to Enlighten} (2nd ed.). Analytics Press.

\bibitem{lundberg2017} Lundberg, S. M., \& Lee, S. I. (2017). A unified approach to interpreting model predictions. In \textit{Advances in Neural Information Processing Systems 30} (NIPS 2017).

\bibitem{ribeiro2016} Ribeiro, M. T., Singh, S., \& Guestrin, C. (2016). ``Why should I trust you?'': Explaining the predictions of any classifier. In \textit{Proceedings of KDD 2016} (pp. 1135--1144).

\bibitem{anscombe1973} Anscombe, F. J. (1973). Graphs in statistical analysis. \textit{The American Statistician}, 27(1), 17--21.

\end{thebibliography}

\end{document}
